% Canonical Paper II source.
The analytic object considered in this section is the three-dimensional contact
state space of Paper I.  It is not a regular AES surface.  We use $D_u,D_v$ for
its horizontal lifts and reserve $X_u,X_v$ for the canonical frame of a
two-dimensional regular AES.  No identification of these frames or of their
associated Laplacians is made here.

Assume throughout that $\mu,\lambda\in\R$ and $\mu\lambda\ne0$, and put
\[
 \mathcal C=\R^3_{(u,v,a)},\qquad
 \alpha=da-\mu\,du-\lambda a\,dv,
\]
\begin{equation}\label{eq:contact-horizontal-data}
 \mathcal D=\ker\alpha,\qquad
 D_u=\partial_u+\mu\partial_a,\qquad
 D_v=\partial_v+\lambda a\partial_a,\qquad T=\partial_a.
\end{equation}
The conventions agree with Paper I, in particular
\begin{equation}\label{eq:contact-bracket-import}
 [D_u,D_v]=\mu\lambda T.
\end{equation}

\subsection{Horizontal metric, complex structure, and Frobenius obstruction}

The contact form alone does not select the analytic normalization below.  We make
the additional choice explicit.

\begin{definition}[Normalized arithmetic contact data]
\label{def:normalized-contact-analytic-data}
The normalized arithmetic contact data consist of $\mathcal D$, the horizontal
metric $h$, and the endomorphism $J$ defined by
\begin{equation}\label{eq:contact-h-and-J}
 h(D_i,D_j)=\delta_{ij},\qquad JD_u=D_v,\qquad JD_v=-D_u.
\end{equation}
The ordered frame $(D_u,D_v)$ fixes the horizontal orientation.  The associated
positive contact measure is
\begin{equation}\label{eq:positive-contact-measure}
 d\nu_{\mathrm c}:=|\alpha\wedge d\alpha|
 =|\mu\lambda|\,du\,dv\,da,
\end{equation}
where the last expression denotes a positive density.
\end{definition}

Thus $J$ is compatible with the chosen horizontal metric and orientation.  We do
not claim that it is selected by the contact distribution without this
normalization, or that it is invariant under arbitrary contactomorphisms.

\begin{proposition}[Frobenius nonintegrability]
\label{prop:contact-frobenius-obstruction}
The distribution $\mathcal D$ is bracket generating of step two and is not
Frobenius integrable.  At every $p\in\mathcal C$,
\[
 \mathcal D_p+[\,\mathcal D,\mathcal D\,]_p=T_p\mathcal C,\qquad
 \alpha([D_u,D_v])=\mu\lambda\ne0.
\]
Consequently no immersed surface in $\mathcal C$ can have tangent bundle equal to
$\mathcal D$ on a nonempty open set.
\end{proposition}

\begin{proof}
The horizontal fields are independent and project to
$\partial_u,\partial_v$.  Equation~\eqref{eq:contact-bracket-import} supplies the
vertical field $T$, so $D_u,D_v,[D_u,D_v]$ span every tangent space.  Since
$\alpha(T)=1$, the bracket is not horizontal.  The Frobenius theorem gives the
remaining assertions.
\end{proof}

This is the sharp reason that the contact branch cannot be regarded as another
coordinate presentation of the surface branch in
Sections~\ref{sec:analytic-data}--\ref{sec:surface-operators}.

\subsection{Adjoints and the variational sub-Laplacian}

Formal-adjoint statements are initially made on
$C_c^\infty(\mathcal C)\subset L^2(\mathcal C,d\nu_{\mathrm c})$, with the inner
product complex linear in its first entry.

\begin{proposition}[Contact-volume divergences and adjoints]
\label{prop:contact-divergence-adjoints}
With respect to $d\nu_{\mathrm c}$,
\begin{equation}\label{eq:contact-divergences}
 \diver_{\nu_{\mathrm c}}D_u=0,\qquad
 \diver_{\nu_{\mathrm c}}D_v=\lambda.
\end{equation}
Consequently,
\begin{equation}\label{eq:contact-first-order-adjoints}
 D_u^*=-D_u,\qquad D_v^*=-D_v-\lambda.
\end{equation}
\end{proposition}

\begin{proof}
The density in \eqref{eq:positive-contact-measure} is a nonzero constant multiple
of Lebesgue measure.  The coordinate divergence of
$D_u=\partial_u+\mu\partial_a$ is zero, whereas that of
$D_v=\partial_v+\lambda a\partial_a$ is
$\partial_a(\lambda a)=\lambda$.  Integration by parts gives
$Y^*=-Y-\diver Y$ for a real vector field $Y$.
\end{proof}

\begin{definition}[Contact energy and variational sub-Laplacian]
\label{def:contact-variational-sublaplacian}
For $f,g\in C_c^\infty(\mathcal C)$, set
\begin{equation}\label{eq:contact-dirichlet-form}
 \mathcal E_{\mathrm c}(f,g)
 =\int_{\mathcal C}
 \bigl(D_uf\,\overline{D_ug}+D_vf\,\overline{D_vg}\bigr)\,d\nu_{\mathrm c}.
\end{equation}
The divergence-of-horizontal-gradient operator and its nonnegative energy operator
are
\begin{align}
 \Delta_{\mathcal C}
   &:=D_u^2+D_v^2+\lambda D_v,\label{eq:contact-variational-laplacian}\\
 A_{\mathcal C}
   &:=-\Delta_{\mathcal C}
     =D_u^*D_u+D_v^*D_v.\label{eq:contact-energy-operator}
\end{align}
\end{definition}

\begin{proposition}[Energy identity and Friedrichs realization]
\label{prop:contact-friedrichs-realization}
The form $\mathcal E_{\mathrm c}$ is densely defined, nonnegative, and closable in
$L^2(\mathcal C,d\nu_{\mathrm c})$.  Its closure determines a unique nonnegative
self-adjoint Friedrichs operator $A_{\mathcal C}^{\mathrm F}$.  On
$C_c^\infty(\mathcal C)$ its differential expression is $A_{\mathcal C}$, and
\begin{equation}\label{eq:contact-energy-identity}
 \langle A_{\mathcal C}f,f\rangle
 =\mathcal E_{\mathrm c}(f,f)
 =\|D_uf\|_2^2+\|D_vf\|_2^2.
\end{equation}
\end{proposition}

\begin{proof}
Define
\[
 \mathsf D:C_c^\infty(\mathcal C)\longrightarrow L^2(d\nu_{\mathrm c})^2,
 \qquad \mathsf Df=(D_uf,D_vf).
\]
Proposition~\ref{prop:contact-divergence-adjoints} shows that
$\Dom(\mathsf D^*)$ contains the dense subspace
$C_c^\infty(\mathcal C)^2$.  Hence $\mathsf D$ is closable, and so is the
quadratic form $\|\mathsf Df\|_2^2$.  The representation theorem for closed
nonnegative forms gives the Friedrichs operator.  Integration by parts gives
\[
 \mathcal E_{\mathrm c}(f,g)
 =\langle(D_u^*D_u+D_v^*D_v)f,g\rangle
 =\langle A_{\mathcal C}f,g\rangle,
\]
which proves the differential expression and the energy identity.
\end{proof}

\begin{theorem}[H\"ormander hypoellipticity]
\label{thm:contact-hormander-hypoellipticity}
The differential expressions $A_{\mathcal C}$ and
$\Delta_{\mathcal C}$ are hypoelliptic.  If $A_{\mathcal C}f$ is smooth on
an open set in the distributional sense, then $f$ is smooth there.
\end{theorem}

\begin{proof}
Proposition~\ref{prop:contact-frobenius-obstruction} verifies H\"ormander's bracket
condition for $D_u,D_v$.  The second-order part is their sum of squares and the
remaining term is first order, so the sum-of-squares theorem applies
\cite{Hormander1967Hypoelliptic}.
\end{proof}

\begin{remark}
This is a local regularity statement.  It does not assert essential
self-adjointness on $C_c^\infty(\mathcal C)$, identify the Friedrichs spectrum,
or solve a boundary problem on an arbitrary contact domain.
\end{remark}

\subsection{Raw and adjoint CR factorizations}

\begin{definition}[Contact CR operators]\label{def:contact-cr-operators}
Set
\begin{equation}\label{eq:contact-cr-definition}
 \dbarC=\frac12(D_u+iD_v),\qquad
 \dC=\frac12(D_u-iD_v),\qquad
 Q_{\mathcal C}=D_u^2+D_v^2.
\end{equation}
A smooth complex field $F$ is contact arithmetic CR when $\dbarC F=0$.
\end{definition}

\begin{proposition}[Raw curvature-twisted factorization]
\label{prop:contact-raw-cr-factorization}
On smooth complex fields,
\begin{align}
 4\dC\dbarC
   &=Q_{\mathcal C}+i\mu\lambda T,\label{eq:contact-raw-factorization-plus}\\
 4\dbarC\dC
   &=Q_{\mathcal C}-i\mu\lambda T.\label{eq:contact-raw-factorization-minus}
\end{align}
Every contact arithmetic CR field therefore satisfies
\begin{equation}\label{eq:contact-twisted-equation}
 \bigl(\Delta_{\mathcal C}-\lambda D_v+i\mu\lambda T\bigr)F=0.
\end{equation}
\end{proposition}

\begin{proof}
Expansion gives
\[
 4\dC\dbarC=(D_u-iD_v)(D_u+iD_v)
 =D_u^2+D_v^2+i[D_u,D_v].
\]
This proves the first identity; reversing the order proves the second.  For a CR
field, apply $\dC$ and use
$Q_{\mathcal C}=\Delta_{\mathcal C}-\lambda D_v$.
\end{proof}

\begin{proposition}[Contact-volume adjoint factorizations]
\label{prop:contact-adjoint-cr-factorization}
As formal identities on
$C_c^\infty(\mathcal C)\subset L^2(d\nu_{\mathrm c})$,
\begin{equation}\label{eq:contact-cr-adjoints}
 \dbarC^*=\frac12(-D_u+iD_v+i\lambda),\qquad
 \dC^*=\frac12(-D_u-iD_v-i\lambda),
\end{equation}
and
\begin{align}
 4\dbarC^*\dbarC
   &=A_{\mathcal C}+i\lambda D_u-i\mu\lambda T,\label{eq:contact-adjoint-factorization-dbar}\\
 4\dC^*\dC
   &=A_{\mathcal C}-i\lambda D_u+i\mu\lambda T.\label{eq:contact-adjoint-factorization-d}
\end{align}
The Reeb field of $\alpha$ is
\begin{equation}\label{eq:contact-Reeb-field}
 \Reeb=-\mu^{-1}\partial_u,
\end{equation}
and the same identities therefore take the intrinsic contact form
\begin{align}
 4\dbarC^*\dbarC
   &=A_{\mathcal C}-i\mu\lambda\Reeb,
   \label{eq:contact-adjoint-factorization-Reeb-dbar}\\
 4\dC^*\dC
   &=A_{\mathcal C}+i\mu\lambda\Reeb.
   \label{eq:contact-adjoint-factorization-Reeb-d}
\end{align}
\end{proposition}

\begin{proof}
The adjoints follow from \eqref{eq:contact-first-order-adjoints}, with scalar
coefficients conjugated.  Expanding the first product,
\begin{align*}
4\dbarC^*\dbarC
 &=(-D_u+iD_v+i\lambda)(D_u+iD_v)\\
 &=-D_u^2-D_v^2-\lambda D_v+i\lambda D_u-i[D_u,D_v]\\
 &=A_{\mathcal C}+i\lambda D_u-i\mu\lambda T.
\end{align*}
The second calculation is analogous and gives the stated signs.
Finally, $d\alpha=-\lambda\,da\wedge dv$ has kernel spanned by
$\partial_u$, and $\alpha(-\mu^{-1}\partial_u)=1$; hence
$\Reeb=-\mu^{-1}\partial_u$.  Substitution of
$D_u=\partial_u+\mu T$ into the two preceding formulas gives their Reeb forms.
\end{proof}

The raw identities display vertical curvature.  The adjoint identities also contain
the measure-dependent drift required by the contact energy.  In particular, every
smooth contact arithmetic CR field satisfies the pointwise equation
\begin{equation}\label{eq:contact-Reeb-twisted-equation}
 (A_{\mathcal C}-i\mu\lambda\Reeb)F=0,
 \qquad\text{equivalently}\qquad
 (\Delta_{\mathcal C}+i\mu\lambda\Reeb)F=0.
\end{equation}
This differential consequence does not assert that $F$ belongs to the
$L^2$-domain of the Friedrichs operator.

\subsection{The balanced measure and its unitary gauge}

\begin{definition}[Balanced horizontal measure]\label{def:contact-balanced-measure}
Define
\begin{equation}\label{eq:contact-balanced-measure}
 d\nu_{\mathrm b}=e^{-\lambda v}\,d\nu_{\mathrm c}.
\end{equation}
Multiplication by
\begin{equation}\label{eq:contact-measure-unitary}
 Uf=e^{-\lambda v/2}f
\end{equation}
defines a unitary map from $L^2(d\nu_{\mathrm b})$ to
$L^2(d\nu_{\mathrm c})$.
\end{definition}

\begin{proposition}[Two-measure gauge identities]
\label{prop:contact-two-measure-gauge}
With respect to $d\nu_{\mathrm b}$, both horizontal fields are divergence free.
Thus $D_u^{*,\mathrm b}=-D_u$, $D_v^{*,\mathrm b}=-D_v$, and
$-Q_{\mathcal C}$ is the balanced-measure energy expression.  Moreover,
\begin{align}
 U D_uU^{-1}&=D_u,\label{eq:gauge-Du}\\
 U D_vU^{-1}&=D_v+\frac{\lambda}{2},\label{eq:gauge-Dv}\\
 U(-Q_{\mathcal C})U^{-1}
   &=A_{\mathcal C}-\frac{\lambda^2}{4},\label{eq:gauge-sum-squares}\\
 U\dbarC U^{-1}&=\dbarC+\frac{i\lambda}{4},\qquad
 U\dC U^{-1}=\dC-\frac{i\lambda}{4}.\label{eq:gauge-cr}
\end{align}
\end{proposition}

\begin{proof}
For $e^\phi d\nu_{\mathrm c}$,
$\diver_{e^\phi\nu_{\mathrm c}}Y=\diver_{\nu_{\mathrm c}}Y+Y\phi$.
With $\phi=-\lambda v$, one has $D_u\phi=0$ and $D_v\phi=-\lambda$,
which cancels \eqref{eq:contact-divergences}.  Unitarity follows directly from
\[
 \int|Uf|^2d\nu_{\mathrm c}
 =\int|f|^2e^{-\lambda v}d\nu_{\mathrm c}.
\]
Since $D_uv=0$ and $D_vv=1$, differentiating
$U^{-1}f=e^{\lambda v/2}f$ proves the first two conjugation formulas.  Hence
\[
 U(-Q_{\mathcal C})U^{-1}
 =-D_u^2-\left(D_v+\frac{\lambda}{2}\right)^2
 =A_{\mathcal C}-\frac{\lambda^2}{4}.
\]
The CR formulas follow from their definitions.
\end{proof}

\begin{warning}
The balanced measure is an arithmetic-coordinate gauge, not the positive contact
measure.  Every statement about adjoints, spectra, or $L^2$-membership must name
the measure used.
\end{warning}

\subsection{Logarithmic CR first integrals and modes}

Choose a holomorphic branch $\mathrm{Log}_{\mu}$ on
\[
 \Omega_\mu=\{w\in\C:\sgn(\mu)\Re w>0\}.
\]
The line $\mu+i\lambda\R$ lies in this half-plane.

\begin{proposition}[Two CR first integrals]
\label{prop:contact-logarithmic-first-integrals}
The fields
\begin{equation}\label{eq:contact-first-integrals}
 z=u+iv,\qquad
 \zeta=u+\frac{i}{\lambda}\mathrm{Log}_{\mu}(\mu+i\lambda a)
\end{equation}
satisfy $\dbarC z=\dbarC\zeta=0$.  The map
\begin{equation}\label{eq:contact-cr-embedding}
 \Psi:\mathcal C\longrightarrow\C^2,\qquad \Psi=(z,\zeta),
\end{equation}
has real rank three everywhere.
\end{proposition}

\begin{proof}
One has $D_uz=1$, $D_vz=i$, and hence
$\dbarC z=\tfrac12(1+i^2)=0$.  Further,
\[
 \partial_a\zeta=-\frac1{\mu+i\lambda a},\quad
 D_u\zeta=\frac{i\lambda a}{\mu+i\lambda a},\quad
 D_v\zeta=-\frac{\lambda a}{\mu+i\lambda a},
\]
so $D_u\zeta+iD_v\zeta=0$.  Finally,
\[
 dz=du+i\,dv,\qquad d\zeta=du-\frac{da}{\mu+i\lambda a}.
\]
If a real tangent vector is killed by both differentials, the first forces its
$u,v$ components to vanish and the second forces its $a$ component to vanish.
Thus $d\Psi$ is injective.
\end{proof}

\begin{corollary}[Explicit contact CR modes]\label{cor:contact-cr-modes}
If $\Phi$ is holomorphic on a neighborhood of $\Psi(\mathcal C)$, then
$\Phi(z,\zeta)$ is contact arithmetic CR.  In particular, for
$\rho,\sigma\in\C$,
\begin{equation}\label{eq:contact-exponential-modes}
 F_{\rho,\sigma}
 =e^{\rho z+\sigma\zeta}
 =e^{(\rho+\sigma)u+i\rho v}
   (\mu+i\lambda a)^{\,i\sigma/\lambda}
\end{equation}
is a smooth CR field in the chosen logarithmic branch and depends genuinely on
$a$ when $\sigma\ne0$.
\end{corollary}

\begin{proof}
The chain rule and Proposition~\ref{prop:contact-logarithmic-first-integrals} give
\[
 \dbarC\Phi(z,\zeta)
 =\Phi_1\dbarC z+\Phi_2\dbarC\zeta=0.
\]
The displayed modes use $\Phi(w_1,w_2)=e^{\rho w_1+\sigma w_2}$.
\end{proof}

No square-integrability, orthogonality, spanning, or completeness assertion is made
for these modes.

\subsection{The filtered affine--Appell module}

Let $r=e^{\lambda v}$, set
\begin{equation}\label{eq:contact-appell-generators}
 B_n=r^{-n}a^n,\qquad n\ge0,
\end{equation}
and let $\mathcal A=C^\infty(\R^2_{(u,v)},\C)$, pulled back to
$\mathcal C$.  Define
\begin{equation}\label{eq:contact-filtered-modules}
 \mathcal M_N
 =\left\{\sum_{n=0}^NP_n(u,v)B_n:P_n\in\mathcal A\right\}.
\end{equation}
We also set $\mathcal M_{-1}=\{0\}$.

\begin{proposition}[Filtered affine--Appell stability]
\label{prop:contact-filtered-appell-module}
The modules $\mathcal M_N$ form an increasing filtration and
$\mathcal M_N\mathcal M_M\subseteq\mathcal M_{N+M}$.  For
$P\in\mathcal A$, with $B_{-1}=0$,
\begin{align}
 D_u(PB_n)
  &=(\partial_uP)B_n+\mu nr^{-1}PB_{n-1},\label{eq:contact-module-Du}\\
 D_v(PB_n)
  &=(\partial_vP)B_n,\label{eq:contact-module-Dv}\\
 T(PB_n)&=nr^{-1}PB_{n-1}.\label{eq:contact-module-T}
\end{align}
Each $\mathcal M_N$ is stable under
$D_u,D_v,\dbarC,\dC,Q_{\mathcal C},\Delta_{\mathcal C}$, and
$A_{\mathcal C}$, while $T\mathcal M_N\subseteq\mathcal M_{N-1}$.
\end{proposition}

\begin{proof}
The product assertion follows from $B_nB_m=B_{n+m}$.  Direct differentiation
gives
\[
 D_uB_n=\mu nr^{-1}B_{n-1},\qquad D_vB_n=0,\qquad
 TB_n=nr^{-1}B_{n-1}.
\]
Since $r^{-1}$ is smooth in $v$, the coefficients remain in $\mathcal A$.
First-order stability follows by linear combination and second-order stability by
iteration.
\end{proof}

\begin{remark}
The term \emph{module} is essential.  The finite-dimensional complex span of
$B_0,\ldots,B_N$ is not stable under $D_u$, because $r^{-1}$ appears in
\eqref{eq:contact-module-Du}.  No density or completeness claim is made.  The
finite upward-sweep antiderivative is proved in
Appendix~\ref{app:contact-calculations}.
\end{remark}
